\SGASection{6}{Parabolic subgroups and split tori}
\SGARefTarget{sga3:xxvi:sec:6}

\begin{sgastatement}{\SGAInlineRefTarget{sga3:xxvi:proposition:6-1}Proposition 6.1}
Let \(S\) be a scheme, \(G\) an \(S\)-reductive group,
\(\mathfrak g=\operatorname{Lie}(G/S)\), and \(Q\) a split subtorus
of \(G\).  Write \(Q=D_S(M)\), and let
\[
 \mathfrak g=\bigoplus_{\alpha\in M}\mathfrak g^\alpha
\]
be the decomposition of \(\mathfrak g\) under the action of \(Q\).
Let \(M_1\subset M\) satisfy
\[
 0\in M_1,\qquad
 \alpha,\beta\in M_1\Longrightarrow\alpha+\beta\in M_1.
\]
\begin{enumerate}[label=\textup{(\roman*)}]
\item There is a unique smooth subgroup \(H_{M_1}\) of \(G\), with
connected fibers, which contains \(\Centr_G(Q)\) and whose Lie
algebra is
\[
 \bigoplus_{\alpha\in M_1}\mathfrak g^\alpha.
\]
\item The following implications hold:
\[
\begin{aligned}
 M_1=\{0\}
   &\Longrightarrow H_{M_1}=\Centr_G(Q),\\
 M_1=-M_1
   &\Longrightarrow H_{M_1}\text{ is reductive},\\
 M_1\cup(-M_1)=M
   &\Longrightarrow
   \begin{gathered}
    H_{M_1}\text{ and }H_{-M_1}\text{ are opposite parabolic}\\
    \text{subgroups of \(G\), with common Levi subgroup }
    H_{M_1\cap(-M_1)}.
   \end{gathered}
\end{aligned}
\]
\end{enumerate}
\end{sgastatement}
\SGARefTarget{sga3:xxvi:prop:6:1}

Assertions \textup{(i)} and \textup{(ii)} are local for the fpqc
topology.  We may therefore suppose that \(Q\) is contained in a
maximal torus \(T\) of \(G\), and may moreover split \(G\) relative
to \(T\).  Assertion \textup{(i)} then follows immediately from
\SGARefLink{sga3:xxii:cor:5.3.5}{Exposé XXII, 5.3.5}, \SGARefLink{sga3:xxii:prop:5.4.5}{5.4.5}, and \SGARefLink{sga3:xxii:thm:5.4.7}{5.4.7}.  The assertions in
\textup{(ii)} follow from \SGARefLink{sga3:xxii:cor:5.3.5}{Exposé XXII, 5.3.5}, \SGARefLink{sga3:xxii:prop:5.10.1}{5.10.1}, and \SGARefLink{sga3:xxii:lem:5.11.3}{5.11.3},
together with \SGARefLink{sga3:xxvi:prop:1:3}{1.4} and \SGARefLink{sga3:xxvi:item:4:3:1}{4.3.2} of the present Exposé.

\begin{sgastatement}{Corollary 6.2}
Let \(S\) be a scheme, \(G\) an \(S\)-reductive group, and \(Q\) a
split subtorus of \(G\).  There is a parabolic subgroup of \(G\)
having \(\Centr_G(Q)\) as a Levi subgroup.
\end{sgastatement}
\SGARefTarget{sga3:xxvi:cor:6:2}

Indeed, write \(Q=D_S(M)\), choose a total order on the group \(M\),
and let \(M_1\) be the set of its positive elements.  The subgroup
\(H_{M_1}\) has the required property.

\begin{sgastatement}{Corollary 6.3}
If the \(S\)-reductive group \(G\) has a noncentral split subtorus,
then it has a proper parabolic subgroup \textup{(that is, one
different from \(G\))}.
\end{sgastatement}
\SGARefTarget{sga3:xxvi:cor:6:3}

Combining \SGARefLink{sga3:xxvi:cor:6:2}{6.2} with \SGARefLink{sga3:xxvi:corollary:5-9}{5.9} and \SGARefLink{sga3:xxvi:corollary:5-10}{5.10} gives the following result.

\begin{sgastatement}{\SGAInlineRefTarget{sga3:xxvi:corollary:6-4}Corollary 6.4}
Let \(S\) be a scheme, \(G\) an \(S\)-reductive group, and \(Q\) a
split subtorus of \(G\).  The canonical morphism
\[
 G\longrightarrow G/\Centr_G(Q)
\]
is a locally trivial fibration.

If \(S\) is semilocal, the map
\[
 G(S)\longrightarrow\bigl(G/\Centr_G(Q)\bigr)(S)
\]
is surjective, and the map
\[
 H^1\bigl(S,\Centr_G(Q)\bigr)\longrightarrow H^1(S,G)
\]
is injective.
\end{sgastatement}
\SGARefTarget{sga3:xxvi:cor:6:4}

\SGASubsection{6.5}{}
\SGARefTarget{sga3:xxvi:subsec:6:5}

Suppose that \(S\) is connected.  If \(T\) is an \(S\)-torus and
\(T'\), \(T''\) are two split subtori of \(T\), their product
\[
 T'\cdot T''
\]
is again a split subtorus of \(T\).\footnote{\emph{Editor's note.}
The Polo--Gille edition adopts multiplicative notation here,
replacing the source's ``sum \(T'+T''\)'' by ``product
\(T'\cdot T''\)''.}  Indeed, it is identified with the quotient of
\(T'\times_S T''\) by \(T'\cap T''\), and this quotient is split by
\SGARefLink{sga3:ix:proposition:2:11}{Exposé IX, 2.11}.  It follows that \(T\) has a largest split subtorus,
which we denote by \(T_{\mathrm{spl}}\).

\begin{sgastatement}{\SGAInlineRefTarget{sga3:xxvi:lemma:6-6}Lemma 6.6}
Let \(S\) be a connected scheme, \(T\) an isotrivial \(S\)-torus,
and \(T_{\mathrm{spl}}\) its largest split subtorus.  The following
conditions are equivalent:
\begin{enumerate}[label=\textup{(\roman*)}]
\item There is a homomorphism
\[
 T\longrightarrow\mathbf G_{m,S}
\]
different from the trivial homomorphism.
\item \(T_{\mathrm{spl}}\ne e\).
\end{enumerate}
\end{sgastatement}
\SGARefTarget{sga3:xxvi:lem:6:6}

Since \(T\) is isotrivial, there are a finite group \(\Gamma\), a
connected principal Galois covering \(S'\to S\) with group
\(\Gamma\), and an isomorphism
\[
 T_{S'}\simeq D_{S'}(M).
\]
The group \(M\) is then a \(\Gamma\)-module, and there is a natural
isomorphism
\[
 \operatorname{Hom}_S(T,\mathbf G_{m,S})
   =H^0(\Gamma,M).
\]

Let \(V\) be the vector subspace of \(M\otimes\mathbf Q\) generated
by the elements
\[
 \gamma(m)-m,\qquad \gamma\in\Gamma,\quad m\in M.
\]
One checks immediately that
\[
 (T_{\mathrm{spl}})_{S'}
   \simeq D_{S'}\bigl(M/(M\cap V)\bigr).
\]
Thus \textup{(i)} is equivalent to
\[
 H^0(\Gamma,M)\ne0,
\]
or equivalently to
\[
 H^0(\Gamma,M\otimes\mathbf Q)\ne0,
\]
whereas \textup{(ii)} is equivalent to \(M\ne M\cap V\), or to
\(M\otimes\mathbf Q\ne V\).  But
\[
 M\otimes\mathbf Q
 =H^0(\Gamma,M\otimes\mathbf Q)\oplus V.
\]
This follows at once by considering the projector from
\(M\otimes\mathbf Q\) onto its invariants which sends \(x\) to the
average of its transforms under \(\Gamma\).

\begin{sgastatement}{Lemma 6.7}
Let \(S\) be a scheme, \(G\) an \(S\)-reductive group, \(P\ne G\) a
parabolic subgroup of \(G\), \(L\) a Levi subgroup of \(P\), and
\(Q=\rad(L)\).  There is a homomorphism
\[
 Q\longrightarrow\mathbf G_{m,S}
\]
different from the trivial homomorphism.
\end{sgastatement}
\SGARefTarget{sga3:xxvi:lem:6:7}

Let \(U=\rad^u(P)\).  It is invariant under \(\Int(P)\), hence under
\(\Int(Q)\).  Consider the invertible \(\mathcal O_S\)-module
\[
 \det\bigl(\operatorname{Lie}(U)\bigr),
\]
the highest exterior power of the locally free
\(\mathcal O_S\)-module \(\operatorname{Lie}(U)\).  The adjoint
representation defines a homomorphism
\[
 f:Q\longrightarrow
 \operatorname{Aut}\!\left(\det\bigl(\operatorname{Lie}(U)\bigr)\right)
 =\mathbf G_{m,S}.
\]
Since \(P\ne G\), one has \(U\ne e\).  Choose \(s\in S\) such that
\(U_s\ne e\).  After splitting \(G_s\) relative to a maximal torus
containing \(Q_s\), it is immediate that \(f_s\ne e\).

\begin{sgastatement}{Proposition 6.8}
Let \(S\) be a connected semilocal scheme, \(G\) an \(S\)-reductive
group, \(P\) a parabolic subgroup of \(G\), \(L\) a Levi subgroup of
\(P\), and \(Q=\rad(L)\).  Let \(Q_{\mathrm{spl}}\) be the largest
split subtorus of \(Q\), equivalently the largest central split
subtorus of \(L\).  Then
\[
 L=\Centr_G(Q_{\mathrm{spl}}).
\]
\end{sgastatement}
\SGARefTarget{sga3:xxvi:prop:6:8}

Put \(L'=\Centr_G(Q_{\mathrm{spl}})\).  This is a reductive subgroup
of \(G\) containing \(L\), and
\[
 P'=P\cap L'
\]
is a parabolic subgroup of \(L'\) with Levi subgroup \(L\), by
\SGARefLink{sga3:xxvi:proposition:1-20}{1.20}.  If \(L'\ne L\), then \(L'\not\subset P\), since \(L\) is a
maximal reductive subgroup of \(P\) by \SGARefLink{sga3:xxvi:prop:1:6}{1.7}.  Hence \(P'\ne L'\).

Let \(G_1\) be the derived group of \(L'\), and put
\(P_1=P'\cap G_1\).  By \SGARefLink{sga3:xxvi:proposition:1-19}{1.19}, \(P_1\) is a parabolic subgroup of the
semisimple group \(G_1\), the group \(L_1=L\cap G_1\) is a Levi
subgroup of \(P_1\), and
\[
 Q_1=\rad(L_1)
    =\bigl(\rad(L)\cap G_1\bigr)^0
    =\bigl(Q\cap G_1\bigr)^0.
\]
The group \(L_1\) has a maximal torus \(T_1\), by \SGARefLink{sga3:xiv:3.20}{Exposé XIV, 3.20},
and \(T_1\) is isotrivial, by \SGARefLink{sga3:xxiv:theorem:4-1-5}{Exposé XXIV, 4.1.5}.  Therefore its
subtorus \(Q_1\) is also isotrivial, by \SGARefLink{sga3:ix:proposition:2:11}{Exposé IX, 2.11}.  Since
\(P_1\ne G_1\), \SGARefLink{sga3:xxvi:lem:6:7}{Lemmas 6.7} and \SGARefLink{sga3:xxvi:lemma:6-6}{6.6} give
\[
 (Q_1)_{\mathrm{spl}}\ne e.
\]
It follows that
\[
 \bigl(Q_{\mathrm{spl}}\cap G_1\bigr)^0\ne e,
\]
and hence that \(Q_{\mathrm{spl}}\not\subset\rad(L')\), because
\(\rad(L')\cap G_1\) is finite.  This contradicts the definition of
\(L'\).
\SGAQED
